\documentclass[10pt,a4paper]{amsart}
\usepackage{amsmath,amssymb,amsthm,graphicx}
\usepackage{verbatim}
\newenvironment{reference}{\comment}{\endcomment}
%\newenvironment{reference}{}{}
\newenvironment{slogan}{\comment}{\endcomment}
\newenvironment{history}{\comment}{\endcomment}
% For commutative diagrams we use Xy-pic
\usepackage[all]{xy}

% We use 2cell for 2-commutative diagrams.
\xyoption{2cell}
% Theorem environments.
%
\theoremstyle{plain}
\newtheorem{theorem}[subsection]{Theorem}
\newtheorem{proposition}[subsection]{Proposition}
\newtheorem{lemma}[subsection]{Lemma}

\theoremstyle{definition}
\newtheorem{definition}[subsection]{Definition}
\newtheorem{example}[subsection]{Example}
\newtheorem{exercise}[subsection]{Exercise}
\newtheorem{situation}[subsection]{Situation}

\theoremstyle{remark}
\newtheorem{remark}[subsection]{Remark}
\newtheorem{remarks}[subsection]{Remarks}

\numberwithin{equation}{subsection}

% Macros
%
\def\lim{\mathop{\mathrm{lim}}\nolimits}
\def\colim{\mathop{\mathrm{colim}}\nolimits}
\def\Spec{\mathop{\mathrm{Spec}}}
\def\Hom{\mathop{\mathrm{Hom}}\nolimits}
\def\Ext{\mathop{\mathrm{Ext}}\nolimits}
\def\SheafHom{\mathop{\mathcal{H}\!\mathit{om}}\nolimits}
\def\SheafExt{\mathop{\mathcal{E}\!\mathit{xt}}\nolimits}
\def\Sch{\mathit{Sch}}
\def\Mor{\mathop{\mathrm{Mor}}\nolimits}
\def\Ob{\mathop{\mathrm{Ob}}\nolimits}
\def\Sh{\mathop{\mathit{Sh}}\nolimits}
\def\NL{\mathop{N\!L}\nolimits}
\def\CH{\mathop{\mathrm{CH}}\nolimits}
\def\proetale{{pro\text{-}\acute{e}tale}}
\def\etale{{\acute{e}tale}}
\def\QCoh{\mathit{QCoh}}
\def\Ker{\mathop{\mathrm{Ker}}}
\def\Im{\mathop{\mathrm{Im}}}
\def\Coker{\mathop{\mathrm{Coker}}}
\def\Coim{\mathop{\mathrm{Coim}}}

% Boxtimes
%
\DeclareMathSymbol{\boxtimes}{\mathbin}{AMSa}{"02}

%
% Macros for moduli stacks/spaces
%
\def\QCohstack{\mathcal{QC}\!\mathit{oh}}
\def\Cohstack{\mathcal{C}\!\mathit{oh}}
\def\Spacesstack{\mathcal{S}\!\mathit{paces}}
\def\Quotfunctor{\mathrm{Quot}}
\def\Hilbfunctor{\mathrm{Hilb}}
\def\Curvesstack{\mathcal{C}\!\mathit{urves}}
\def\Polarizedstack{\mathcal{P}\!\mathit{olarized}}
\def\Complexesstack{\mathcal{C}\!\mathit{omplexes}}
% \Pic is the operator that assigns to X its picard group, usage \Pic(X)
% \Picardstack_{X/B} denotes the Picard stack of X over B
% \Picardfunctor_{X/B} denotes the Picard functor of X over B
\def\Pic{\mathop{\mathrm{Pic}}\nolimits}
\def\Picardstack{\mathcal{P}\!\mathit{ic}}
\def\Picardfunctor{\mathrm{Pic}}
\def\Deformationcategory{\mathcal{D}\!\mathit{ef}}
\usepackage{cleveref}
\crefname{originallemma}{original lemma}{original lemmas}
\Crefname{originallemma}{Original lemma}{Original lemmas}
\makeatletter
\newlabel{divisors-lemma-characterize-OD}{{01X0}{source}{}{sourcepointer-1}{}}
\newlabel{divisors-lemma-characterize-OD@cref}{{[originallemma][0][]01X0}{[1][0][]source}{}{}{}}
\makeatother
\crefname{originallemma}{original lemma}{original lemmas}
\Crefname{originallemma}{Original lemma}{Original lemmas}
\makeatletter
\newlabel{lemma-chow-ring-projective-bundle}{{02TX}{source}{}{sourcepointer-2}{}}
\newlabel{lemma-chow-ring-projective-bundle@cref}{{[originallemma][0][]02TX}{[1][0][]source}{}{}{}}
\makeatother
\crefname{originallemma}{original lemma}{original lemmas}
\Crefname{originallemma}{Original lemma}{Original lemmas}
\makeatletter
\newlabel{lemma-linebundle}{{0B74}{source}{}{sourcepointer-3}{}}
\newlabel{lemma-linebundle@cref}{{[originallemma][0][]0B74}{[1][0][]source}{}{}{}}
\makeatother
\crefname{originallemma}{original lemma}{original lemmas}
\Crefname{originallemma}{Original lemma}{Original lemmas}
\makeatletter
\newlabel{lemma-pullback-affine-fibres-surjective}{{02TT}{source}{}{sourcepointer-4}{}}
\newlabel{lemma-pullback-affine-fibres-surjective@cref}{{[originallemma][0][]02TT}{[1][0][]source}{}{}{}}
\makeatother
\crefname{originallemma}{original lemma}{original lemmas}
\Crefname{originallemma}{Original lemma}{Original lemmas}
\makeatletter
\newlabel{lemma-pushforward-cap-c1}{{02SU}{source}{}{sourcepointer-5}{}}
\newlabel{lemma-pushforward-cap-c1@cref}{{[originallemma][0][]02SU}{[1][0][]source}{}{}{}}
\makeatother
\crefname{originallemma}{original lemma}{original lemmas}
\Crefname{originallemma}{Original lemma}{Original lemmas}
\makeatletter
\newlabel{lemma-relative-effective-cartier}{{02TR}{source}{}{sourcepointer-6}{}}
\newlabel{lemma-relative-effective-cartier@cref}{{[originallemma][0][]02TR}{[1][0][]source}{}{}{}}
\makeatother
\crefname{originallemma}{original lemma}{original lemmas}
\Crefname{originallemma}{Original lemma}{Original lemmas}
\makeatletter
\newlabel{lemma-restrict-to-open}{{02RX}{source}{}{sourcepointer-7}{}}
\newlabel{lemma-restrict-to-open@cref}{{[originallemma][0][]02RX}{[1][0][]source}{}{}{}}
\makeatother
\crefname{originallemma}{original lemma}{original lemmas}
\Crefname{originallemma}{Original lemma}{Original lemmas}
\makeatletter
\newlabel{situation-setup}{{02QL}{source}{}{sourcepointer-8}{}}
\newlabel{situation-setup@cref}{{[originallemma][0][]02QL}{[1][0][]source}{}{}{}}
\makeatother
\title{Flat pullback along a vector bundle is an isomorphism on Chow groups}
\author{The Stacks Project Authors}
\date{stacks-project:chow.tex:local-snapshot-20261012; GNU FDL 1.2 or later}
\begin{document}
\maketitle
\noindent\textit{Editorial note.} All delivered mathematical text is the author's own, copied from the pinned Stacks Project chapter file chow.tex; no statement or proof step was rewritten, reordered or removed.\par\medskip
\noindent\textit{Editorial note.} Preamble provenance: the pinned chapter file opens with \textbackslash{}input\{preamble\} and carries no preamble text of its own, so the project's single Stacks-wide preamble staged at sources/stacks-chow/preamble.tex is excerpted verbatim as three spans, lines 8-12 (the verbatim package together with the project's own reference, slogan and history display-environment definitions, including the author's commented-out older form of the reference environment), lines 14-18 (Xy-pic setup) and lines 154-217 (the native theorem declarations and mathematical macros). The lines 8-12 span is required because a counted statement or its surrounding text may use those display environments, which the project typesets as comments. Nothing was added, removed or redefined; the wrapper class and the amsmath/amssymb/amsthm/graphicx base packages are supplied by the assembler.\par\medskip
\noindent\textit{Editorial note.} Same-chapter pointers are located at their \textbackslash{}label lines inside the pinned chapter file; cross-chapter pointers are located in the project's own tag list staged in the snapshot (sources/stacks-chow/tags.txt), at the verbatim lines carrying their canonical tags.\par\medskip
\noindent\textit{Allowed standard prerequisites (editorial).}\begin{itemize}
\item Surjectivity of pullback along a linear fibration, applied here to the line bundle case. Same chapter chow.tex, labels lemma-linebundle (proof line 7078) and lemma-pullback-affine-fibres-surjective (proof line 7079). Gives surjectivity of the pullback map p\textasciicircum{}* directly, since a vector bundle is a flat morphism with affine space fibres; the injectivity computation is what the counted span supplies. The affine fibres lemma is claimed separately as S14-algebra-pullback-affine-fibres.
\item The relation between the divisor of a section and the twisting sheaf, and restriction to an open. Divisors, label divisors-lemma-characterize-OD (proof line 7092); Same chapter chow.tex, label lemma-restrict-to-open (proof line 7104). Identifies O\_P(D) with O\_P(1) for the effective Cartier divisor cut out by the section, and describes the kernel of the restriction of Chow groups to the complement of D as the cycles supported on D.
\item The projective space bundle formula, the description of the Gysin pullback, and the vanishing of the top cap on the bundle. Same chapter chow.tex, labels lemma-chow-ring-projective-bundle (proof line 7106), lemma-relative-effective-cartier (proof line 7112) and lemma-pushforward-cap-c1 (proof lines 7113 and 7122). Writes a cycle on D through the powers of the first Chern class of the restriction of the twisting sheaf, converts that expression into powers on P, and then observes that the top power cap vanishes on the projective bundle of rank one higher, forcing all the classes to be zero. The bundle formula is claimed separately as S14-algebra-projective-bundle-formula.
\end{itemize}\medskip
\ifdefined\hypertarget\hypertarget{sourcepointer-1}{}\fi
\noindent\textit{Source reference.} Original lemma 01X0: sources/stacks-chow/tags.txt, source lines 2501--2501. Navigation pointer for \textbackslash{}ref\{divisors-lemma-characterize-OD\} used inside the counted span (line 7092); the argument that consumes it is retained in the span.\par\medskip
\ifdefined\hypertarget\hypertarget{sourcepointer-2}{}\fi
\noindent\textit{Source reference.} Original lemma 02TX: sources/stacks-chow/pinned-original-source.tex, source lines 6932--6932. Navigation pointer for \textbackslash{}ref\{lemma-chow-ring-projective-bundle\} used inside the counted span (line 7106); the argument that consumes it is retained in the span.\par\medskip
\ifdefined\hypertarget\hypertarget{sourcepointer-3}{}\fi
\noindent\textit{Source reference.} Original lemma 0B74: sources/stacks-chow/pinned-original-source.tex, source lines 5737--5737. Navigation pointer for \textbackslash{}ref\{lemma-linebundle\} used inside the counted span (line 7078); the argument that consumes it is retained in the span.\par\medskip
\ifdefined\hypertarget\hypertarget{sourcepointer-4}{}\fi
\noindent\textit{Source reference.} Original lemma 02TT: sources/stacks-chow/pinned-original-source.tex, source lines 5651--5651. Navigation pointer for \textbackslash{}ref\{lemma-pullback-affine-fibres-surjective\} used inside the counted span (line 7079); the argument that consumes it is retained in the span.\par\medskip
\ifdefined\hypertarget\hypertarget{sourcepointer-5}{}\fi
\noindent\textit{Source reference.} Original lemma 02SU: sources/stacks-chow/pinned-original-source.tex, source lines 4631--4631. Navigation pointer for \textbackslash{}ref\{lemma-pushforward-cap-c1\} used inside the counted span (line 7113); the argument that consumes it is retained in the span.\par\medskip
\ifdefined\hypertarget\hypertarget{sourcepointer-6}{}\fi
\noindent\textit{Source reference.} Original lemma 02TR: sources/stacks-chow/pinned-original-source.tex, source lines 5598--5598. Navigation pointer for \textbackslash{}ref\{lemma-relative-effective-cartier\} used inside the counted span (line 7112); the argument that consumes it is retained in the span.\par\medskip
\ifdefined\hypertarget\hypertarget{sourcepointer-7}{}\fi
\noindent\textit{Source reference.} Original lemma 02RX: sources/stacks-chow/pinned-original-source.tex, source lines 2949--2949. Navigation pointer for \textbackslash{}ref\{lemma-restrict-to-open\} used inside the counted span (line 7104); the argument that consumes it is retained in the span.\par\medskip
\ifdefined\hypertarget\hypertarget{sourcepointer-8}{}\fi
\noindent\textit{Source reference.} Original lemma 02QL: sources/stacks-chow/pinned-original-source.tex, source lines 1372--1372. Navigation pointer for \textbackslash{}ref\{situation-setup\} used inside the counted span (line 7063); the argument that consumes it is retained in the span.\par\medskip
\begin{lemma}
\label{lemma-vectorbundle}
Let $(S, \delta)$ be as in Situation \ref{situation-setup}.
Let $X$ be locally of finite type over $S$.
Let $\mathcal{E}$ be a finite locally free sheaf of rank $r$ on $X$.
Let
$$
p :
E = \underline{\Spec}(\text{Sym}^*(\mathcal{E}))
\longrightarrow
X
$$
be the associated vector bundle over $X$.
Then $p^* : \CH_k(X) \to \CH_{k + r}(E)$ is an isomorphism for all $k$.
\end{lemma}

\begin{proof}
(For the case of linebundles, see Lemma \ref{lemma-linebundle}.)
For surjectivity see Lemma \ref{lemma-pullback-affine-fibres-surjective}.
Let $(\pi  : P \to X, \mathcal{O}_P(1))$
be the projective space bundle associated
to the finite locally free sheaf $\mathcal{E} \oplus \mathcal{O}_X$.
Let $s \in \Gamma(P, \mathcal{O}_P(1))$ correspond to the global
section $(0, 1) \in \Gamma(X, \mathcal{E} \oplus \mathcal{O}_X)$.
Let $D = Z(s) \subset P$. Note that
$(\pi|_D : D \to X , \mathcal{O}_P(1)|_D)$
is the projective space bundle associated
to $\mathcal{E}$. We denote $\pi_D = \pi|_D$ and
$\mathcal{O}_D(1) = \mathcal{O}_P(1)|_D$.
Moreover, $D$ is an effective
Cartier divisor on $P$. Hence $\mathcal{O}_P(D) = \mathcal{O}_P(1)$
(see Divisors, Lemma \ref{divisors-lemma-characterize-OD}).
Also there is an isomorphism
$E \cong P \setminus D$. Denote $j : E \to P$ the
corresponding open immersion.
For injectivity we use that the kernel of
$$
j^* :
\CH_{k + r}(P)
\longrightarrow
\CH_{k + r}(E)
$$
are the cycles supported in the effective Cartier divisor $D$,
see Lemma \ref{lemma-restrict-to-open}. So if $p^*\alpha = 0$, then
$\pi^*\alpha = i_*\beta$ for some $\beta \in \CH_{k + r}(D)$.
By Lemma \ref{lemma-chow-ring-projective-bundle} we may write
$$
\beta = \pi_D^*\beta_0 +
\ldots + c_1(\mathcal{O}_D(1))^{r - 1} \cap \pi_D^* \beta_{r - 1}.
$$
for some $\beta_i \in \CH_{k + i}(X)$.
By Lemmas \ref{lemma-relative-effective-cartier}
and \ref{lemma-pushforward-cap-c1}
this implies
$$
\pi^*\alpha = i_*\beta =
c_1(\mathcal{O}_P(1)) \cap \pi^*\beta_0 +
\ldots +
c_1(\mathcal{O}_D(1))^r \cap \pi^*\beta_{r - 1}.
$$
Since the rank of $\mathcal{E} \oplus \mathcal{O}_X$ is $r + 1$
this contradicts Lemma \ref{lemma-pushforward-cap-c1} unless all
$\alpha$ and all $\beta_i$ are zero.
\end{proof}

\end{document}
